% 2007 年福建卷文科第 22 题：定点 F(1,0)、直线 l:x=-1，P 为动点，Q 为 P 在 l 上的
%   垂足，且 QP·QF = FP·FQ；第（2）问过 F 的直线交轨迹 C 于 A、B，交 l 于 M。
% 来源 PDF：gaokaomath-src/普通高考/2007/2007福建文.pdf
% SHA-256：96ffeae7fae0016108ed33e565abdf5983e06aef19f6452a8b4c8c5e5e4e58d0
% 原页：2007福建文.pdf 第 2 页（A3 三栏）第 3 栏，书页码 370。
% 重绘依据：按原页只画坐标轴、准线 l 与焦点 F，并在 x 轴上标出 -1 与 1。
\begin{tikzpicture}[x=0.85cm, y=0.85cm, semithick,
  every node/.style={inner sep=1.5pt}]

  % 坐标轴
  \draw[->,>=stealth] (-3.2,0) -- (4.2,0) node[right] {$x$};
  \draw[->,>=stealth] (0,-2.4) -- (0,3.0) node[above] {$y$};
  \node[below left] at (0,0) {$O$};

  % 直线 l：x=-1
  \draw (-1,-2.0) -- (-1,2.6) node[above] {$l$};

  % 焦点 F
  \fill (1,0) circle (1.3pt);
  \node[above right=1pt] at (1,0) {$F$};

  % x 轴上的刻度 -1、1
  \node[below] at (-1,0) {$-1$};
  \node[below] at (1,0) {$1$};
\end{tikzpicture}
